\subsection*{先改变付款路径，再解释成功}

先看一笔独立的购买：购买库里还没有购买记录，Alice 在支付平台的付款账户有 CNY 10000 分。购买库保存账户 1 的业务归属和购买资料，这笔购买的资金变化则全部记在平台账本中。平台扣去 3000 分后，付款账户余下 7000 分，购买服务不再同时扣减本地账户余额。\footnote{本章采用一次全额付款模型。平台完成收款的判定随所用支付方式确定，授权冻结、捕获、清算与结算可分别属于不同阶段。本章模拟器把成功限定为一次已提交的扣款流水。}平台付款账户通过已经验证的付款关联，对应到购买服务的账户 1。

付款方式变化以后，接口需要区分受理购买与完成付款。前六章的 \code{purchase.create.v1} 在成功终态中同时报告本地订单、扣款与原结果。本章为外部支付引入动作 \code{purchase.\allowbreak create.\allowbreak external.v2}，结果按 \code{purchase-result-external-v2} 解释。完整身份仍为
\[
 S=(\text{account\_id},\text{action},\text{session\_id},\text{idempotency\_key}),
\]
这次购买取 \(S=(1,\text{purchase.create.external.v2},\text{E1},\text{K1})\)。

E1 仍是有固定期限的登记会话，K1 在这次确认的第一次发送以前持久保存。规范化输入继续采用 \code{purchase-input-v1} 的约定，包含商品、数量 1、币种 CNY 和报价 Q7。经验证的报价确定金额为 3000 分，本地受理事务同时保存原输入与价格快照，后来的支付和核对都引用这份金额。

外部支付接口 v2 只处理用户新确认的购买。已经通过 v1 发起的购买，即使结果尚未查明，也沿原来的 \(S\) 和 v1 入口恢复。把动作名改成 v2，会使完整身份中的 \code{action} 改变，旧购买仍按原来的处理路径继续。假如客户端保留 E1、K1，却把同一笔未决购买改发给 v2，服务就会把它识别为另一项购买：旧路径可能完成本地扣款，新路径也可能完成平台收款。图~\ref{fig:ch07-version} 把这个变化放在完整身份中比较。

\CSevenFigure{ch07-version}{动作名属于操作身份，旧购买沿原入口恢复}{fig:ch07-version}

受理入口继续核对已登记会话、操作许可、账户和动作，准入与会话持久关闭协调。有效恢复先查同一 \(S\) 并比较旧输入，新 \(S\) 才进入首次受理。关闭或缺失 E1 时，购买入口按第三章规则拒绝请求。对于此前已经受理的 C1，服务端恢复和经授权的业务事实查询继续读取原映射与证据，核查它的付款结局。

本地首次受理时，受约束插入取得 \(S\)，建立 P1，并将 \code{business\_state} 设为 \code{awaiting\_payment}。这里的 \code{business\_state} 表示购买在本地的业务阶段，\code{awaiting\_payment} 表示购买已经受理，付款尚未取得足够的终态证据。受理事务同时分配支付操作 C1，它标识 P1 需要完成的这一次全额收款，与 P1 一对一关联，供发送、回调、查询和对账共同定位。

本章采用付款确认后建立订单的业务规则。本地受理时只有 P1、C1，还没有订单 O1。取得足够的成功证据以后，本地在保存支付结论的事务中建立 O1，把它关联到 P1、C1 和平台收款 X1，同时保存支付结果，把 P1 改为 \code{paid}。O1 的初始履约状态为 \code{unfulfilled}，后续履约再推进它。付款最终拒绝只保存拒绝结论，不建立 O1。资金流水归平台账本，本地没有第一部分那条余额扣款 D1。

\subsection*{三次提交怎样分布在两个系统}

将本地受理提交记为 A，平台资金提交记为 E，本地保存支付结论的提交记为 R。一条正常历史按 A、E、R 推进。A 共同保存 \(S\)、P1、C1、原输入及下游身份映射。发送者读取已提交的 C1 后调用平台，E 使平台付款流水和余额变化成立。R 则共同保存平台证据、唯一订单 O1 与本地支付结果，使 P1 进入 \code{paid}。

\CSevenFigure{ch07-original01}{本地受理、平台扣款和本地保存结论分别提交}{fig:ch07-commits}

两个提交域，是指各自能够独立决定提交或回滚的一组数据。购买库拥有一个提交域，平台账本拥有另一个。普通本地事务的回滚只撤销购买库内的修改。即使把平台调用写在 \code{BEGIN} 和 \code{COMMIT} 之间，E 成功以后，本地连接仍可能断开并回滚。平台没有参加共同提交协议时，本地数据库没有替平台撤销 E 的步骤。

展开这条历史：购买库先提交 A，P1 等待付款。工作进程调用平台，E 把余额从 10000 改为 7000，建立平台收款 X1。工作进程随后在本地事务内填写成功，尚未提交就退出。新连接读取购买库，P1 仍等待付款；独立连接读取平台库，X1 和 3000 分流水已经存在。那次尚未提交的本地更新随连接结束而回滚，平台流水保持原样。

两边分别持久保存自己的提交。恢复者重新连接时，能够读到本地等待与平台成功并存。重新打开本地事务也不会把两份历史合并为一次共同提交，因此设计要容许这段间隔，并为恢复者保留定位原收款的依据。

第一部分把 \code{pending} 限定在未提交的事务内，因为那笔事务能够同时决定购买、扣款和原结果。这里若把 P1、C1 也留在未提交事务里，E 以后进程退出，原支付映射就会随本地回滚丢失，恢复者甚至不知道该查哪次收款。因此先提交 A，固定已经受理的购买及支付身份，再等待外部结局。\code{awaiting\_payment} 成为可恢复的持久业务阶段，既容许平台尚未执行，也容许 E 已完成而 R 尚未提交。

这个中间状态有明确的迁移规则：A 声明购买已经受理；R 根据可信证据保存订单和结论；证据不足时保留原身份。持久记录保持可读，恢复路径重新发现未决 C1，才能继续推进。第八章将展开重新发现和执行资格，本章先确定每次恢复能够采取哪些动作。

若先把本地写成 \code{paid} 再调用平台，则会留下相反的历史：本地已经报告付款成功，E 却尚未发生，平台余额仍为 10000。先提交等待状态，才能准确表达此时已经成立的受理事实。

\subsection*{恢复者面对的两条历史}

让服务在 A 以后只读本地 P1，并查看上次请求超时。历史 \(H_0\) 中，支付请求到达平台以前就丢失，平台余额仍为 10000；历史 \(H_1\) 中，平台已提交 X1，余额为 7000，返回消息丢失。两条历史都在 R 以前结束，购买库都只有相同的 P1、C1 和等待状态，服务都没有完整平台响应。

\CSevenFigure{ch07-original02}{相同的本地等待状态可以对应未扣款和已扣款}{fig:ch07-indistinguishable}

将此时恢复者能够读取的资料记为观察投影 \(\mathcal V(H)\)，其中包含原输入、\(S\)、C1、本地状态和这次超时。两条历史满足
\[
 \mathcal V(H_0)=\mathcal V(H_1).
\]
假如平台只有普通的可重复执行扣款接口，立即再扣可能使 \(H_1\) 留下两笔扣款；永久停止发送又会让 \(H_0\) 始终没有付款。只要没有新证据，也没有终止原请求未来执行资格的机制，等待更久仍可能得到相同观察。

Birrell 与 Nelson 对通信异常的分析，同样保留了调用可能执行一次或尚未执行的分支~\cite{rpc}。这里又多了一层独立持久状态：平台过程被调用、平台提交扣款、本地保存结论，分别由不同事件决定。恢复者可以寻找区分这些事件的证据，也可以利用平台对同一操作的仲裁，让两条历史采用同一种安全的继续动作。

A 确认以后，接口可以返回 202，在响应体中给出 P1、C1 和 \code{awaiting\_payment}。随后查询 P1，即使得到 200，响应体也可能仍在等待付款；这次查询只是成功读到了记录。等到服务核对付款成功的证据并提交 R，响应体才报告 \code{paid}。此后履约服务根据已确认付款的订单发货或授予商品访问权，并在订单中另行记录履约进度。
